\chapter{Systemzugang}\label{systemzugang}

Das Frontend ist ohne lokale Installation über folgende URL erreichbar:

\begin{center}
\url{https://liveticker-frontend.onrender.com}
\end{center}

\noindent Die KI-Textgenerierung setzt eine aktive n8n-Workflow-Instanz voraus. Eine vollständige Live-Demo ist auf Anfrage beim Autor verfügbar.

\chapter[Interviewprotokoll]{Interviewprotokoll: Experteninterview Senior Manager Redaktion und Digital Media bei Eintracht Frankfurt}\label{interviewleitfaden-experteninterview-sportredakteur-eintracht-frankfurt}

Geführt am 22.~Mai~2026 mit dem Senior Manager Redaktion und Digital Media bei Eintracht Frankfurt. Strukturiertes Leitfadeninterview mit offenen Fragen, qualitative Validierung der in Kapitel~2 hergeleiteten Anforderungen sowie Bewertung des Systems aus Sicht des Praxispartners (Abschnitt~\ref{experteninterview-auswertung}). Antworten kursiv.

\begin{center}\rule{0.5\linewidth}{0.5pt}\end{center}

\section{Einstieg}

\textbf{IF0:} Wie stehen Sie generell zum Einsatz von KI im Sportjournalismus -- sehen Sie es eher als Unterstützung oder als Risiko für die redaktionelle Qualität?

\textit{Unterstützend ist KI hilfreich, die redaktionelle Qualitätssicherung aller Inhalte sollte aber weiterhin komplett in den Händen der Redakteur:innen liegen. Zumal KI bei der Erstellung der Inhalte durch die angesteuerten Token mehr auf Wahrscheinlichkeit als auf Wahrheit setzt, sprich die Fehlerquote insbesondere in einem auch von Zahlen getriebenen Feld wie dem Profifußball zu hoch ist. KI ist hilfreich, zum Beispiel als automatisierter und umfangreicher Datenlieferant bei Live-Ereignissen im Hintergrund, kann aber menschlichen Input -- vor allem hinsichtlich Emotionen und dem jeweiligen (Spannungs-)umfeld eines Klubs angepassten Inhalt und Sprachgebrauch -- nicht ersetzen.}

\begin{center}\rule{0.5\linewidth}{0.5pt}\end{center}

\section{Block 1: Redaktioneller Workflow und Aufwand}

\textbf{IF1:} Wie läuft die Liveticker-Produktion bei Eintracht Frankfurt aktuell ab? Beschreiben Sie bitte einen typischen Spieltag.

\textit{Der Liveticker wird etwa zwei Tage vor dem Matchday mit ersten kleinen Einträgen gestartet. Am Matchday nimmt der Ticker dann schon vor Anpfiff Fahrt auf: Statistiken, Serien, Hintergrundinfos, Impressionen aus dem Stadion kurz vor Anpfiff und allgemein Lust auf das Spiel machende Inhalte. Mit Anpfiff geht's dann ans Kerngeschäft.}

\textbf{IF1a:} Wie bereiten Sie sich vor einem Spieltag vor? Gibt es ein festes Briefing, und wie lange vor Anpfiff beginnen Sie mit der Vorbereitung?

\textit{Die Vorbereitung auf das nächste Spiel zieht sich über die gesamte Woche durch zahlreiche weitere Inhalte, die auf die Website bzw.\ in die App einlaufen. Ein festes Briefing gibt es nicht, die Informationsquellen sind jedoch meist gleich. Bei Englischen Wochen verkürzt sich das entsprechende Zeitfenster.}

\textbf{IF1b:} Von wo aus wird typischerweise getickert -- aus dem Stadion, vom Heimarbeitsplatz, aus dem Hotel bei Auswärtsspielen? Macht das einen Unterschied für den Workflow?

\textit{Getickert wird immer im Stadion, auch bei Auswärtsspielen. Direkt im Stadion bekommt man die meisten Eindrücke, hat alles im Blick, saugt die Stimmung auf und ist zudem nicht auf Live-Übertragungen und -Streams angewiesen. Als in der Saison 2024/25 das Auswärtsspiel in Bochum mit über einer Stunde Verspätung angepfiffen wurde, haben wir die höchsten Zugriffszahlen auf den Liveticker in der gesamten Saison erzielt.}

\textbf{IF2:} Wie viele Personen sind an einem Spieltag für den Liveticker eingeteilt?

\textit{Zwei: Ein Redakteur, der tickert, und ein technischer Administrator, der sicherstellt, dass alles reibungslos funktioniert. Unser Website-Dienstleister stellt zusätzlich während der Spiele einen technischen Support zur Verfügung.}

\textbf{IF3:} Welche Aufgaben muss ein Ticker-Autor neben dem Tickern parallel erledigen?

\textit{Parallel übernimmt der Ticker-Autor keine weiteren redaktionellen Aufgaben, der volle Fokus liegt während des Spiels auf dem Liveticker. Weitere Aufgaben stehen nur vor und nach dem Spiel an, doch nicht während des Spiels.}

\textbf{IF3a:} Wie viele Ticker-Einträge erstellen Sie typischerweise pro Spiel? Gibt es eine Zielzahl -- oder richtet sich das rein nach der Ereignisdichte?

\textit{Das lässt sich konkret nicht festmachen, da es auch stark vom Spielverlauf und der Spielgeschichte abhängt. Zudem hat jeder Redakteur einen etwas anderen Stil bzw.\ eine andere Herangehensweise.}

\textbf{IF3b:} Wie viel Zeit verbringen Sie im Schnitt mit dem Schreiben eines einzelnen Eintrags -- und wie viel Zeit entfällt auf das reine Beobachten des Spiels? Was fühlt sich stressiger an?

\textit{Ein Eintrag kann zwischen zwei und ca.\ 30 Sekunden in Anspruch nehmen, von einem Wort bis zur detaillierten Schilderung einer Chance oder eines Treffers. Faktor Zeit: abhängig vom Spielverlauf und der Spielgeschichte. Stressig fühlt es sich nicht an, es ist sehr geübt. Stress ist bei einem schnellen Medium wie dem Liveticker ein sehr schlechter Ratgeber.}

\textbf{IF3c:} Wie viel Zeit vergeht typischerweise zwischen einem Ereignis -- Tor, Einwechslung, Rote Karte -- und der Veröffentlichung des zugehörigen Eintrags? Gibt es einen internen Richtwert oder Erwartungsdruck, wie schnell reagiert werden soll?

\textit{Es gibt keinen internen Richtwert. Man ist sehr darauf bedacht, die wichtigsten Ereignisse so zeitnah wie möglich rauszuschießen, die Qualität darf dabei aber nicht leiden. In der Regel sind Ereignisse spätestens zwischen 30 und 90 Sekunden später online -- je nach Wortwahl auch schneller. Geschieht viel parallel, so priorisiert man und zieht gegebenenfalls mit entsprechendem Wording nach.}

\textbf{IF4:} Was sind die häufigsten Fehlerquellen unter Zeitdruck? Passieren zum Beispiel Zahlendreher, falsche Spielernamen oder stilistische Brüche?

\textit{Kleinere Rechtschreibfehler.}

\textbf{IF4a:} Welche Tools und Systeme nutzen Sie aktuell für den Liveticker -- gibt es ein CMS, ein eigenes Tool, oder arbeiten Sie in einem externen System? Und wie würde sich ein neues Werkzeug wie dieses dort einfügen?

\textit{Wir benutzen ein eigenes, gebautes Liveticker-Backend von unserem technischen Dienstleister für die Website. Einträge und Pushes etc.\ laufen über eine Schnittstelle in die App rein.}

\textbf{IF4b:} Was sind die größten Schwachstellen im aktuellen Workflow? Gibt es Medientypen -- Spielstandsicons, Statistiken, Fotos, Social-Media-Posts -- die besonders umständlich einzubinden sind, oder Schritte, die unverhältnismäßig viel Zeit kosten?

\textit{Prinzipiell lässt sich die Benutzeroberfläche sehr intuitiv bedienen, zudem sind wir sehr geübt -- umständlich bzw.\ mit zu großem Zeitaufwand verbunden ist keiner der Arbeitsschritte. Bugs wurden im Laufe der Saison beseitigt. Die größte Schwachstelle: schwache Internetverbindung.}

\textbf{IF4c:} Woher beziehen Sie Spielstandsdaten und Statistiken in Echtzeit -- gibt es eine Datenfeed-Integration, oder werden diese manuell übertragen? Wie zuverlässig ist diese Quelle unter Spieltagsbedingungen?

\textit{Wir nutzen Opta und DFBnet. Die Daten werden in der Regel automatisch übertragen, nur in Einzelfällen bei extremer Kurzfristigkeit -- zum Beispiel Verletzung beim Aufwärmen -- müssen Redakteure Einträge manuell überschreiben und korrigieren.}

\textbf{IF4d:} Wie würden Sie das aktuelle System vom Aussehen beschreiben und wie lauten die aktuellen Pain-Points? Werden dort automatisch Statistiken, Aufstellungen oder andere relevante Sachen aktualisiert?

\textit{Der Liveticker ist in das Matchcenter integriert. In anderen Reitern des Matchcenters laufen Statistiken und Aufstellungen inklusive Wechseln, Karten etc.\ automatisch ein.}

\begin{center}\rule{0.5\linewidth}{0.5pt}\end{center}

\section{Block 2: Mehrsprachigkeit und Stilistik}

\textbf{IF5:} Gibt es aktuell Bedarf an mehrsprachigen Tickern? Wenn ja, welche Sprachen wären relevant?

\textit{Englisch wäre sinnvoll. Problem bei automatisierter Übersetzung durch KI: klassischer hessischer umgangssprachlicher Wortschatz bzw.\ Fußball-typische Begriffe oder künstlich langgezogene Worte wie zum Beispiel \glqq{}Toooooooooooor\grqq{} oder \glqq{}Jaaaaaaaa\grqq{}.}

\textbf{IF6:} Wie wird bisher sichergestellt, dass verschiedene Autoren in einem konsistenten Stil schreiben? Gibt es einen Styleguide? Falls ja, wie lautet dieser (wenn möglich)?

\textit{Einen bestimmten Stilguide gibt es hinsichtlich der Schreibweise einzelner Begriffe, jedoch nicht bezüglich der Wortwahl allgemein. Drei Redakteure teilen sich den Liveticker über eine gesamte Saison auf. Und jeder Redakteur pflegt einen anderen Schreibstil bzw.\ eine andere Herangehensweise, auch die Taktung der Einträge kann sich unterscheiden. Insgesamt ist es aber der Liveticker auf eintracht.de, aus der Übung und Erfahrung heraus bewegen wir uns alle in einem gelernten Gesamtgerüst -- bei aller sportlicher Rivalität basierend auf respektvollen, gesunden Umgangsformen.}

\textbf{IF7:} Wie würden Sie den Eintracht-Ticker-Stil beschreiben: Was macht ihn aus, was unterscheidet ihn von neutraler Berichterstattung?

\textit{Durch die Eintracht-Brille. Der Eintracht-Ticker-Stil ist emotionaler und parteiisch (pro Eintracht), aber keinesfalls unsachlich oder gar beleidigend. Er ist fan-nah, mit klarer Eintracht-Sprache, typisch Hessisch. Information steht über allem, aber wie diese verpackt wird, liegt an jedem Redakteur. Je nach Spielverlauf bleibt auch mal Zeit für Witz, Schabernack oder lustige Randgeschichten.}

\textbf{IF7a:} Schreiben Sie jeden Eintrag vollständig neu, oder greifen Sie auf feste Formulierungsmuster zurück -- zum Beispiel für Torankündigungen, Wechsel oder Spielendmeldungen? Wie stark variiert der Wortlaut von Eintrag zu Eintrag, und ist starke Variation überhaupt erwünscht? Denken Sie sich im Vorfeld schon Formulierungen aus oder kommen alle spontan?

\textit{Jeder Eintrag ist neu. Kein Ticker gleicht dem anderen. Natürlich wiederholen sich Fußball-typische Begriffe, da auch der Pool an Synonymen irgendwann erschöpft ist. Hier kann ich nur für mich sprechen: alles ist spontan.}

\begin{center}\rule{0.5\linewidth}{0.5pt}\end{center}

\section{Block 3: Systembewertung}

\textbf{IF8:} Wäre es Ihrer Meinung nach möglich, die Qualität der KI-generierten Ticker-Texte ähnlich zu manuell verfassten Texten zu bekommen -- mit guten Prompts, für Standardevents, Vorberichterstattungen und Spielphasenänderungen?

\textit{Da könnte ich jetzt ausholen. Ich glaube möglich ist es auf jeden Fall, aber es müssen sehr viele Faktoren stimmen, damit die Qualität am Ende so gut ist, wie von einem Menschen. Darüber hinaus folgt die Frage, ob man das will. Die Kollegen haben Spaß daran, das ist ihre Leidenschaft, und die sollten wir nicht wegnehmen. Und am Ende ist das Schreiben des Tickers eine Aufgabe, die viel Kreativität erfordert, und da sehe ich keine KI. Ich sehe die KI wie gesagt als Unterstützer, der Inspiration gibt, Vorschläge macht, Fehler erkennt\ldots\ Bei Vorberichten kann sie auch unterstützen, aber Live ist dann nochmal etwas ganz anderes. Da muss man ja auch sehr situativ die Einträge formulieren. Oder wenn es überraschende Ereignisse gibt, wie zum Beispiel: Wenn wir 4:0 führen, es plötzlich einen Notarzteinsatz im Stadion gibt und die Einträge trotzdem total euphorisch sind, weil die KI weiß, dass sie bei Führungen so schreiben soll -- was machst du dann? Es gibt Dinge, da sollte der Mensch im Driverseat sein.}

\textbf{IF9:} Lesen sich die Texte wie Eintracht-Ticker oder eher wie generische Sportberichte?

\textit{Eher generisch. Unser Eintracht-Ticker ist sehr individuell. Ich glaube nicht, dass eine KI das ersetzen kann und soll.}

\textbf{IF9a:} Schätzen Sie: Welcher Anteil der Ticker-Einträge basiert auf strukturierten Daten, die ein System automatisch erkennen könnte -- Tore, Karten, Wechsel -- und welcher Anteil erfordert, was nur ein Mensch wahrnehmen kann: Atmosphäre, Zweikampfverhalten, Stimmung auf dem Platz?

\textit{Der ist glaube ich sehr groß. Aber wie gesagt, es ist nicht das Gleiche.}

\textbf{IF10:} Ist der Coop-Modus (also KI-Entwurf mit redaktioneller Freigabe) ein realistischer Kompromiss zwischen Zeitersparnis und Qualitätskontrolle?

\textit{Das auf jeden Fall, weil die KI in diesem Fall als Unterstützer fungiert.}

\textbf{IF10a:} Wie bewerten Sie die Benutzeroberfläche des Systems -- ist das Layout intuitiv, und würden Sie die Steuerung im Live-Betrieb als praktikabel einschätzen?

\textit{Die finde ich sehr gut. Sehr intuitiv und nutzerfreundlich. Aus meiner UX-Sicht war das ziemlich gut umgesetzt.}

\textbf{IF10b:} Wie bewerten Sie das Gesamtkonzept des Systems? Was überzeugt Sie, was würden Sie grundlegend anders lösen?

\textit{Siehe oben. Das sah, was die Usability oder UX generell angeht, sehr gut aus. Unseres ist glaube ich nochmal etwas komplexer. Da gibt es auch parallele Streams und Nutzungen, und ein Rechtesystem brauchen wir auch (wer darf was tickern), damit Redakteure nicht auf jedes Spiel zugreifen können.}

\textbf{IF11:} Welche konkreten Features fehlen grundsätzlich für einen produktiven Einsatz?

\textit{Ich kann mich leider nicht mehr so gut erinnern. Ich habe sehr viele Features gesehen, die wir ebenso nutzen oder gerne hätten. Wir machen uns auch Gedanken über Einbindungen anderer Inhalte, die wir schon haben. Zum Beispiel eine MOTM-Abstimmung oder Partner-Integration. Aber grundlegend war da schon sehr viel dabei. Daher fällt mir nicht wirklich etwas ein.}

\textbf{IF12:} Haben Sie das Gefühl, dass das System aus echten Eintracht-Texten gelernt hat -- erkennen Sie den vereinstypischen Stil wieder, etwa die TOOOOR-Konvention, die SGE-Abkürzung oder die emotionale Tonlage? Oder wirkt es eher wie ein generischer Sportbericht?

\textit{Ich glaube nein. Es fehlt ja auch komplett die Perspektive, die die Redakteure aus dem Stadion haben. Was sie beobachten -- der Trainer gestikuliert gerade wild in Richtung Schiedsrichter, oder die Fans sind gerade alle aufgestanden, weil jetzt Chandler kommt, die Wolken ziehen sich im Moment zusammen usw. -- da fehlt sehr viel Inhalt, den sich die KI nirgendwo ziehen kann, oder wo sie dann anfängt zu halluzinieren.}

\begin{center}\rule{0.5\linewidth}{0.5pt}\end{center}

\section{Block 4: Zukunftsperspektive}

\textbf{IF13:} Könnten Sie sich vorstellen, KI-generierte Ticker im Regelbetrieb zu nutzen?

\textit{Nein. Zumindest nicht im Bereich Profi-Männer und Profi-Frauen. Mit Blick auf Jugendfußball ließe sich zumindest darüber diskutieren.}

\textbf{IF13a:} Sollten KI-generierte Ticker-Einträge für Leser als solche gekennzeichnet sein -- oder spielt die Herkunft des Textes keine Rolle, solange er korrekt und stilistisch passend ist?

\textit{Es käme auf das Ausmaß des Stilbruchs an. Die Herkunft des Textes spielt auf jeden Fall eine Rolle, siehe auch Thema „Wahrscheinlichkeit vs.\ Wahrheit" bei der Erstellung von KI-Inhalten.}

\begin{center}\rule{0.5\linewidth}{0.5pt}\end{center}

\section{Abschluss}

\textbf{IF14:} Haben Sie noch Anmerkungen, Fragen oder Aspekte, die im Interview nicht zur Sprache gekommen sind?

\textit{Nein. Besten Dank.}